% !TeX root = ../papers/nc-p3.tex
\section{特殊な母線を含む三重交点と滑らかな核曲線の決定}
\label{sec:smooth-kernel-curve-exclusion}
\begin{theorem}\label{thm:smooth-kernel-curve-exclusion}
四辺有限の標準有限Frobenius窓で $D$ は整な反対八次式、階数1軌跡は有限とする。
点窓が被約で一般核像 $Z$ が滑らかな曲線なら、$Z$ は有理正規三次曲線である。
従って被約で階数1軌跡が有限な九・十余因子の場合に曲線核があれば、それは必ず特異である。
また前節の $a=0$ の場合、点窓が非被約でも曲線核は有理正規三次であり、
三重母線の入射スキームそのものが点窓になる。
\end{theorem}
\begin{proof}[繊維長三と例外母線]
前節の有限な正規化 $f:S\to D$ と $M\subset f_*L$ を用いる。
まず $a=0$ なら $M=f_*L$ であり、$X$ への制限は局所自由階数3である。
階数一行列の単元成分を消去すればこれはスキーム上でも成立する。
従って $f^{-1}(X)$ は $X$ 上有限平坦次数3である。実際有限局所環の各因子で $L$ を自明化できるので、
線束の押出しの平坦性は構造層のそれと同じである。$S\subset\mathbb P^3\times B$ ゆえこの次数三因子族は
$X\to B^{(3)}$ を定める（滑らかな曲線の長さ三のHilbertスキーム）。
入力を併記すれば $X$ が三重母線入射に閉埋込みされる。前節と同じ有限零スキームの長さ $20-4g$ により
$g=0$ で等号である。$a=0$ から $\delta=3$ となる。この段階では $Z$ の滑らかさも被約な印も要らない。

以下 $a>0$、$B=Z$ は滑らか、点窓は被約とする。前節の双対合成を有限正規化上で読むと
$2L-2H-K_S$ の非零切断である。一般母線での次数は零なので、零因子は有限本の母線からなる。
従って $M\subset f_*L$ の真の商は有限本の入力母線に支えを持つ。商は純一次元である。
両層はどちらも深さ2で（後者は滑らかな曲面からの有限押出し）、商は閉点で深さ1以上となる。
Hilbert多項式 $a(t+1)$ より少なくとも一本が実際の台になる。これらの台の直線を $J_i$ と書く。

$J_i$ の一般点で $f$ の繊維長は3以上である。繊維は $Z$ の部分スキームとして埋め込まれ、
長さ2以下なら核座標の四一次式がその全切断を生成するため中山の補題で商が零となる。
半連続性より $J_i$ の全点でも長さ3以上である。これらの外では長さは $4-\operatorname{rank}N$ なので、
全ての印が長さ三因子に載り、長さ3以上になる入力は $J_i$ と有限な階数1軌跡に限られる。
\end{proof}
\begin{proof}[一つの固定母線と動く母線]
$Z$ がある二次錐に入る場合を先に処理する。その頂点の吹上げ $\mathbb F_2$ で $Z$ の真変換を
$b C_0+c f$ と書けば $\delta=c$, $g=(b-1)(c-b-1)$。真変換はもとの滑らかな曲線と同型である。
頂点を避ければ $c-2b=0$、通れば頂点イデアルが曲線上で単純点のイデアルとなるので $c-2b=1$。
それぞれ $a=-(b-2)^2$ または $a=-(b-1)(b-2)$ であり、いずれも正にならない。

従って $Z$ は二次錐に含まれない。各 $J_i$ に対し、相異なる母線が $J_i$ 上の動点でこれと交わることはない。
もしあれば動く側のパラメータは曲線 $B$ 全体で稠密である。$J_i$ 自身を与える核点を $z_0$ とすると、
$N|_{J_i}$ は $z_0$ を消す。一般階数1は仮定に反し、3なら交点で動く側の核も $z_0$ になり矛盾する。
従って一般階数2で、$z_0$ を除いた三行二次式の余因子は原始多項式列 $b$ の平方の多項式倍である。
$J_i$ 上の二項式で余因子の次数が4なので $\deg b\le2$。
動く側の核の $z_0$ からの射影はこの $[b]$ の像に入る。点または直線なら $Z$ は退化、
非共線な場合は平面二次曲線になり $Z$ が $z_0$ を頂点とする二次錐に入ってしまう。

これより $J_i$ の一般点の繊維は $J_i$ 自身を与える有限個の $B$ の点だけに支えを持つ。
$B$ は滑らかなので、この繊維に含まれる長さ三因子には有限通りしかない。
例外入力も各 $J_i$ で有限である。実際 $J_i$ と交わる異なる母線が無限本あれば交点は動くか固定であり、
固定の場合は全母線が同一点を通ってしまう。各例外入力で $f$ はやはり有限である。
従って三重母線の退化は $B^{(3)}$ の上では有限台である。一方 $J_i$ の一般点が含む三因子の一つは
$J_i$ の全点で共通である（有限個の候補と入射の閉性）。その三因子に対応する入力線型空間はちょうど $J_i$、
三因子の少なくとも一つの台での母線を含有する範囲に制限されるからである。
\end{proof}
\begin{proof}[三次数の余核と全二十点]
$B^{(3)}$ 上の写像 $k^4\otimes\mathcal O\to F^{[3]}$ に対し、転置の六つの線型関係の
内部次数3のKoszul複体を用いる。末端余核 $U$ は転置余核の対称三乗である。
階数4の開部分では核が二元なのでこの次数は完全であり、滑らかな三次元底で非完全台が有限、
自由複体長が3であるため末端余核以外は完全となる。長さは $\int c_3(F^{[3]})=20-4g$。
実際階数六束を $G$ とすればK類は $20-10G^*+4\bigwedge^2 G^*-\bigwedge^3 G^*$ であり、
分裂原理で三次未満のChern指標は零、三次は $c_3(G)$ である。
各印上で繊維に含まれる三因子を一つ選ぶ。各三因子に割り当てた印の三次式への評価により
$U$ はこれら二十個の剰余体の直和へ全射する。実際その閉繊維は対応する入力線型空間上の全三次式で、
入力三次評価が全印に全射だからである。よって $20-4g\ge20$、従って $g=0$ で全射は同型である。
しかし先の $J_i$ の全点で共通な三因子における $U$ の閉繊維は直線上の三次式の四次元で、
割り当てられる印は直線切断上界によって三個以下である（割当てが空でも同じ）。
これは同型に反する。ゆえに $a>0$ はない。
\end{proof}
